% compile_check.tex — Quick visual check of all generated figures and tables
% Compile with: pdflatex compile_check.tex (from the repo root)
% Requires: output/figures/*.png and output/latex/*.tex to exist
\documentclass[11pt]{article}
\usepackage[margin=1in]{geometry}
\usepackage{graphicx}
\usepackage{amsmath}
\usepackage{threeparttable}
\usepackage{booktabs}
\usepackage{multirow}
\usepackage{float}
\usepackage{hyperref}
\usepackage{caption}

\graphicspath{{output/figures/}}

\title{Reproduction Pipeline Output Check\\
\large He (2026) Senior Thesis — Figure and Table Verification}
\author{Auto-generated by \texttt{compile\_check.tex}}
\date{\today}

\begin{document}
\maketitle
\tableofcontents
\clearpage

%%%%%%%%%%%%%%%%%%%%%%%%%%%%%%%%%%%%%%%%%%%%%%%%%%%%%%%%%%%%%%%%%%%%%%%%%%%%%%%
\section{Tables (from \texttt{output/latex/})}
%%%%%%%%%%%%%%%%%%%%%%%%%%%%%%%%%%%%%%%%%%%%%%%%%%%%%%%%%%%%%%%%%%%%%%%%%%%%%%%

\subsection{Table 1: Dataset Summary by Laboratory}
\setcounter{table}{0}  % next table will be Table 1
\input{output/latex/table_lab_summary.tex}

\clearpage
\subsection{Table 5: Chunk Characteristics}
\setcounter{table}{4}  % next table will be Table 5
\input{output/latex/table_chunk_characteristics.tex}

\clearpage
\subsection{Table 6: Within-Session Learning Rates ($\sigma_{\text{trial}}$)}
\setcounter{table}{5}
\input{output/latex/table_sigma_trial.tex}

\clearpage
\subsection{Table 7: Overnight Learning Rates ($\sigma_{\text{day}}$)}
\setcounter{table}{6}
\input{output/latex/table_sigma_day.tex}

\clearpage
\subsection{Table 8: Complete Effect Size Summary}
\setcounter{table}{7}
\input{output/latex/table_effect_sizes.tex}

\clearpage
\subsection{Table 9: Mixed-Effects Model}
\setcounter{table}{8}
\input{output/latex/table_mixed_effects.tex}

\clearpage
\subsection{Table 10: Behavioral Phenotype Distribution}
\setcounter{table}{9}
\input{output/latex/table_phenotype.tex}

\clearpage
\subsection{Table 11: Overnight Jump Magnitudes}
\setcounter{table}{10}
\input{output/latex/table_overnight_magnitude.tex}

\clearpage
\subsection{Table 12: Overnight Jump Direction}
\setcounter{table}{11}
\input{output/latex/table_overnight_direction.tex}

\clearpage
\subsection{Table 14: Phenotype Validation}
\setcounter{table}{13}  % next table will be Table 14
\input{output/latex/table_phenotype_validation.tex}

%%%%%%%%%%%%%%%%%%%%%%%%%%%%%%%%%%%%%%%%%%%%%%%%%%%%%%%%%%%%%%%%%%%%%%%%%%%%%%%
\clearpage
\section{Figures (from \texttt{output/figures/})}
%%%%%%%%%%%%%%%%%%%%%%%%%%%%%%%%%%%%%%%%%%%%%%%%%%%%%%%%%%%%%%%%%%%%%%%%%%%%%%%

\subsection{Figure 2: Dataset Overview}
\begin{figure}[H]
\centering
\includegraphics[width=\textwidth]{fig_dataset_overview.png}
\caption*{\textbf{Dataset Overview.} (A) Selection tiers by laboratory: Tier 1 (all chunks $\geq$1,000 trials), Tier 2 (some 500--1,000), Excluded. (B) Chunking strategy distribution (92.4\% standard / 2.2\% alt / 5.4\% invalid). (C) Trial counts by chunk with tier threshold lines. (D) Total training duration distribution.}
\end{figure}

\clearpage
\subsection{Figure 3: $\sigma$ Trajectories}
\begin{figure}[H]
\centering
\includegraphics[width=\textwidth]{fig_sigma_trajectories_connected.png}
\caption*{\textbf{Learning Rate Parameters Across Training Phases.} Within-session ($\sigma_{\text{trial}}$, top row) and overnight ($\sigma_{\text{day}}$, bottom row) smoothness parameters for each weight across chunks. Bold red lines show population mean $\pm$ SE; gray lines show individual mouse trajectories ($n = 53$). Asterisks indicate significant C1 vs.\ C3 comparisons (Wilcoxon signed-rank, Bonferroni-corrected).}
\end{figure}

\clearpage
\subsection{Figure 4: Trial-by-Trial Weight Dynamics (9 mice)}
\begin{figure}[H]
\centering
\includegraphics[width=\textwidth]{fig_weight_dynamics.png}
\caption*{\textbf{Trial-by-Trial Weight Dynamics.} Weight trajectories across training for 9 mice with balanced chunk durations. Background shading indicates training phase (blue = C1, green = C2, pink = C3); dashed vertical lines mark chunk boundaries. Shaded bands show $\pm$1 SD confidence intervals.}
\end{figure}

\clearpage
\subsection{Figure 5: Contrast Curriculum Progression (9 mice)}
\begin{figure}[H]
\centering
\includegraphics[width=\textwidth]{fig_contrast_curriculum.png}
\caption*{\textbf{Contrast Curriculum Progression.} Daily trial counts by contrast level for the same 9 mice. Stacked bars show trial distribution: 100\% and 50\% dominate C1; 25\% and 12.5\% are introduced in C2; 6.25\% and 0\% appear in C3. Background shading and dashed lines indicate chunk boundaries.}
\end{figure}

\clearpage
\subsection{Figure 6: Individual Composite (3 mice)}
\begin{figure}[H]
\centering
\includegraphics[width=\textwidth]{fig_individual_composite_v3.png}
\caption*{\textbf{Example Animals: Weight Dynamics and Learning Rate Trajectories.} Three mice spanning the range of contrast $\sigma$ trajectories: increasing (top, ZM\_1367), typical/median (middle, ibl\_witten\_07), stable/decreasing (bottom, KS024). \textit{Left}: trial-by-trial weight trajectories with chunk shading and $\pm$1 SD bands. \textit{Center}: $\sigma_{\text{trial}}$ across chunks. \textit{Right}: $\sigma_{\text{day}}$ across chunks.}
\end{figure}

\clearpage
\subsection{Figure 7: Phenotype Distribution}
\begin{figure}[H]
\centering
\includegraphics[width=\textwidth]{fig_phenotype_distribution.png}
\caption*{\textbf{Behavioral Phenotype Distribution.} (A) Distribution of $\log_{10}(\sigma_{\text{pc}}/\sigma_{\text{wsls}})$ ratio across mice ($n = 53$). Dashed vertical lines indicate phenotype thresholds at $\pm 0.48$. (B) Phenotype proportions by chunk.}
\end{figure}

\clearpage
\subsection{Figure 8: Phenotype Stability}
\begin{figure}[H]
\centering
\includegraphics[width=0.7\textwidth]{fig_phenotype_stability.png}
\caption*{\textbf{Phenotype Stability: Chunk 1 vs.\ Chunk 3.} Each point represents one mouse; position indicates $\log_{10}(\sigma_{\text{pc}}/\sigma_{\text{wsls}})$ in C1 (x-axis) and C3 (y-axis). Dotted lines indicate thresholds at $\pm 0.48$. Colors indicate C3 phenotype classification. Points along the diagonal would indicate perfect stability.}
\end{figure}

\clearpage
\subsection{Figure 9: Phenotype \& Learning Speed}
\begin{figure}[H]
\centering
\includegraphics[width=0.7\textwidth]{fig_phenotype_learning_speed.png}
\caption*{\textbf{Phenotype and Learning Speed.} Sessions to reach Chunk 3 by behavioral phenotype. Individual data points overlaid on transparent box plots. Kruskal-Wallis test reported.}
\end{figure}

\clearpage
\subsection{Figure 10: Overnight Analysis}
\begin{figure}[H]
\centering
\includegraphics[width=\textwidth]{fig_overnight_corrected_v3.png}
\caption*{\textbf{Overnight Weight Changes by Training Stage.} Top row: absolute overnight jump magnitudes ($|\Delta\beta|$) by chunk. Bottom row: jump direction using weight-appropriate metrics---contrast: \% positive signed $\Delta W$ (strengthening); bias: \% with $\Delta|\beta| < 0$ (toward zero). Dashed line = 50\% chance; stars indicate pooled binomial significance.}
\end{figure}

\clearpage
\subsection{Figure 11: Overnight Distributions}
\begin{figure}[H]
\centering
\includegraphics[width=\textwidth]{fig_overnight_distributions.png}
\caption*{\textbf{Distributions of Overnight Weight Changes by Training Phase.} Step histograms (density-normalized) of overnight changes for each weight across C1/C2/C3. Light shading marks the ``good'' direction. Inset: median, \% good direction, and sample size per chunk.}
\end{figure}

\clearpage
\subsection{Figure 12: Day Gap Effect}
\begin{figure}[H]
\centering
\includegraphics[width=0.6\textwidth]{fig_day_gap_effect.png}
\caption*{\textbf{Effect of Day Gap on Overnight Jump Magnitude.} Mean absolute weight change as a function of days between sessions (1, 2, 3+). Error bars show SE. Spearman correlations reported as inset.}
\end{figure}

\clearpage
\subsection{Figure 13: Individual Trajectories}
\begin{figure}[H]
\centering
\includegraphics[width=0.8\textwidth]{fig_individual_trajectories.png}
\caption*{\textbf{Individual Contrast Learning Rate Trajectories.} Gray lines show $\sigma_{\text{trial,contrast}}$ for individual mice ($n = 53$); red line shows population mean $\pm$ SE. CV (SD/mean) across mice reported per chunk.}
\end{figure}

\clearpage
\subsection{Figure 14: Lab Effects}
\begin{figure}[H]
\centering
\includegraphics[width=0.8\textwidth]{fig_lab_effects.png}
\caption*{\textbf{Laboratory Effects on Contrast Learning Rate.} $\sigma_{\text{trial,contrast}}$ in Chunk 3 by laboratory (9 labs). Box plots show median and IQR; red points show individual mice. Dashed line indicates grand mean. Kruskal-Wallis test reported.}
\end{figure}

\clearpage
\subsection{Figure 15: Outlier Trajectories}
\begin{figure}[H]
\centering
\includegraphics[width=0.8\textwidth]{fig_outlier_trajectories.png}
\caption*{\textbf{Outlier Identification.} $\sigma_{\text{trial,contrast}}$ trajectories with outliers highlighted (red). Gray lines show non-outlier mice; green connected line shows corrected population mean $\pm$ SE excluding outliers. Orange dotted lines indicate 3 SD thresholds per chunk.}
\end{figure}

\clearpage
\subsection{Figure 16: High-Contrast Control}
\begin{figure}[H]
\centering
\includegraphics[width=\textwidth]{fig_high_contrast_control.png}
\caption*{\textbf{High-Contrast Control.} (A) Paired comparison of $\sigma_{\text{trial,contrast}}$ between C1 and C3-high ($\geq$25\% only); individual data points and medians shown; $W$, $p$, $r$, $n$ reported. (B) Distribution of per-mouse differences ($\Delta\sigma_{\text{trial,contrast}}$, C3-high $-$ C1) with median line. (C) Box plots comparing C1, C3-high, and C3 (original) with transparent boxes and visible data points.}
\end{figure}

\clearpage
\subsection{Figure 25: Phenotype Validation}
\begin{figure}[H]
\centering
\includegraphics[width=\textwidth]{fig_phenotype_validation.png}
\caption*{\textbf{Phenotype Validation: Empirical Win-Stay/Lose-Stay Rates.} (A) Scatter of $\log_{10}$-ratio vs WS$-$LS gap with linear regression line; Spearman $\rho$ and $p$ reported. (B) Grouped box plots by phenotype with transparent boxes, visible data points, and Kruskal-Wallis $H$ and $p$.}
\end{figure}

\end{document}
